% !TEX root = ../main.tex
\begin{ouverture}[mecanisme]{5}{Mécanismes réactionnels}{Partie 1 · Constitution et transformations de la matière}
{Acte élémentaire\enspace·\enspace Intermédiaire réactionnel\enspace·\enspace Flèche courbe\enspace·\enspace Chocs efficaces}
\objectif{Identifier les actes élémentaires, un intermédiaire réactionnel et un catalyseur dans un mécanisme fourni.}
\objectif{Établir l'équation de la réaction modélisée par un mécanisme réactionnel.}
\objectif{Repérer les sites donneurs et accepteurs de doublets d'électrons à l'aide des électronégativités.}
\objectif{Représenter les flèches courbes d'un acte élémentaire en justifiant leur sens.}
\objectif{Interpréter l'influence de la température et des concentrations en termes de chocs efficaces.}
\prerequis{Schémas de Lewis, électronégativité, polarisation d'une liaison.}
\prerequis{Réactions acide-base et couples acide-base (chapitre 1).}
\prerequis{Vitesses volumiques, facteurs cinétiques et catalyse (chapitre 4).}
\end{ouverture}
\essentiel
\essSection{1}{Du bilan macroscopique au mécanisme}
\begin{definition}[Acte élémentaire, mécanisme réactionnel]
L'équation d'une réaction traduit un bilan : elle ne dit pas comment les entités réagissent à l'échelle
microscopique. Un \textbf{acte élémentaire} est une étape qui se produit réellement à cette échelle, lors d'un choc
entre deux entités au plus, ou par la transformation d'une entité unique. Le \textbf{mécanisme réactionnel} est la
succession des actes élémentaires qui modélise la transformation.
\end{definition}
\begin{definition}[Intermédiaire réactionnel]
Un intermédiaire réactionnel est une entité formée au cours d'un acte élémentaire puis consommée au cours d'un acte
suivant. Souvent très réactif et de courte durée de vie, il n'apparaît pas dans l'équation de la réaction.
\end{definition}
\begin{savoirfaire}[Établir l'équation à partir d'un mécanisme]
Additionner les équations des actes élémentaires, en les multipliant au besoin par un même nombre pour que les
intermédiaires se compensent, puis simplifier les entités présentes des deux côtés : les intermédiaires réactionnels
et les catalyseurs disparaissent du bilan.
\end{savoirfaire}
\begin{exemple}
Au-dessous de 500 K, la réaction \ce{NO2 + CO -> NO + CO2} se déroule en deux actes élémentaires :
\[\ce{2 NO2 -> NO3 + NO}\qquad\text{puis}\qquad\ce{NO3 + CO -> NO2 + CO2}.\]
Leur somme, après simplification de \ce{NO3} et d'une molécule de \ce{NO2}, redonne l'équation de la réaction ;
\ce{NO3} est un intermédiaire réactionnel.
\end{exemple}
\essSection{2}{Sites donneurs et accepteurs, flèches courbes}
\begin{definition}[Électronégativité et polarisation]
L'électronégativité $\chi$ d'un élément traduit l'aptitude de ses atomes à attirer les électrons d'une liaison.
Une liaison entre deux atomes d'électronégativités différentes (en pratique $\Delta\chi\geq0{,}4$) est polarisée :
l'atome le plus électronégatif porte une charge partielle négative $\delta^-$, l'autre une charge partielle positive $\delta^+$.
\end{definition}
\schema{05-polarisation}
\begin{propriete}[Sites donneurs et accepteurs de doublets d'électrons]
Un \textbf{site donneur} de doublet d'électrons est un atome porteur d'un doublet non liant ou d'une charge négative
(entière ou partielle), ou une liaison multiple. Un \textbf{site accepteur} est un atome porteur d'une charge positive,
entière ou partielle, comme un atome de carbone lié à un atome plus électronégatif, ou un atome d'hydrogène lié à un
atome d'oxygène, d'azote ou d'halogène.
\end{propriete}
\begin{definition}[Flèche courbe]
Au cours d'un acte élémentaire, des liaisons se forment ou se rompent par déplacement de doublets d'électrons. Une
flèche courbe représente le déplacement d'un doublet : elle part du doublet (non liant ou liant) d'un site donneur et
pointe vers le site accepteur, où se forme la nouvelle liaison, ou vers l'atome qui récupère le doublet d'une liaison
qui se rompt.
\end{definition}
\schema{05-fleches}
\begin{savoirfaire}[Représenter les flèches courbes d'un acte élémentaire]
Repérer, grâce aux électronégativités et aux charges, le site donneur et le site accepteur ; tracer une flèche du doublet
vers l'atome accepteur. Si cet atome a déjà sa couche externe saturée (hydrogène, carbone tétravalent), une de ses
liaisons doit se rompre en même temps : tracer une seconde flèche, de cette liaison vers l'atome qui garde le doublet.
Vérifier enfin que la charge totale est la même avant et après l'acte.
\end{savoirfaire}
\begin{attention}
Une flèche courbe part toujours d'un doublet d'électrons et pointe vers un atome : elle ne part jamais d'un site
accepteur, par exemple d'un atome d'hydrogène. Elle ne doit pas être confondue avec la flèche d'une équation de réaction.
\end{attention}
\essSection{3}{Catalyse et mécanisme}
\begin{propriete}[Rôle microscopique d'un catalyseur]
Un catalyseur modifie le mécanisme réactionnel : il ouvre un autre chemin, formé d'actes élémentaires plus rapides.
Consommé au cours d'un acte élémentaire, il est régénéré au cours d'un acte suivant : il n'apparaît donc pas dans
l'équation de la réaction.
\end{propriete}
\begin{exemple}
La dismutation du peroxyde d'hydrogène, très lente, est catalysée par les ions iodure : ils sont consommés par
l'acte 1 et régénérés par l'acte 2, tandis que l'ion hypoiodite \ce{IO-} est un intermédiaire réactionnel.
\end{exemple}
\schema{05-cycle}
\begin{savoirfaire}[Distinguer catalyseur et intermédiaire réactionnel]
Un catalyseur est présent dès le début : il est consommé dans un acte élémentaire, puis reformé dans un acte suivant.
Un intermédiaire est d'abord formé, puis consommé : il est absent au début comme à la fin de la transformation.
\end{savoirfaire}
\essSection{4}{Interprétation microscopique des facteurs cinétiques}
\begin{propriete}[Chocs efficaces]
Un acte élémentaire n'a lieu que lors d'un choc entre les entités réactives, et seule une petite fraction des chocs
est \textbf{efficace} : il faut une énergie suffisante et une orientation favorable. La vitesse d'un acte élémentaire
augmente avec la \textbf{fréquence des chocs} et avec la \textbf{proportion de chocs efficaces}.
\end{propriete}
\begin{propriete}[Concentration et température]
Quand la concentration des réactifs augmente, les entités sont plus nombreuses par unité de volume : les chocs sont
plus fréquents. Quand la température augmente, les entités se déplacent plus vite : les chocs sont un peu plus
fréquents et, surtout, plus énergétiques, donc plus souvent efficaces.
\end{propriete}
\schema{05-chocs}
\begin{rediger}
Pour justifier une flèche courbe, écrire par exemple : « l'atome d'oxygène, plus électronégatif que l'atome de
carbone, porte des doublets non liants : c'est un site donneur ; la flèche part de l'un de ces doublets et pointe vers
l'atome de carbone, site accepteur ($\delta^+$) ». Pour un facteur cinétique, relier toujours la vitesse aux chocs entre
entités : leur fréquence et leur efficacité.
\end{rediger}
\exercices

\rubrique{Découverte}
\begin{blocLivre}{ex-5-1}
\begin{exercice}[1]{Lire un mécanisme}
Au-dessous de 500 K, la réaction entre le dioxyde d'azote et le monoxyde de carbone est modélisée par le mécanisme :
\[\text{acte 1 : }\ce{NO2 + NO2 -> NO3 + NO}\qquad\qquad\text{acte 2 : }\ce{NO3 + CO -> NO2 + CO2}\]
\begin{questions}
\item Combien d'entités se rencontrent au cours de chaque acte élémentaire ?
\item Identifier l'intermédiaire réactionnel. Justifier.
\item Établir l'équation de la réaction.
\item Le dioxyde d'azote est formé au cours de l'acte 2. Est-il pour autant un catalyseur ? Justifier.
\end{questions}
\end{exercice}
\end{blocLivre}
\begin{blocLivre}{ex-5-2}
\begin{exercice}[1]{Repérer des charges partielles}
On considère le chlorométhane \ce{CH3-Cl}, le méthanol \ce{CH3-OH}, le bromure d'hydrogène \ce{H-Br} et la
méthylamine \ce{CH3-NH2}.
\begin{donnees}
Électronégativités : H : 2,20 ; C : 2,55 ; N : 3,04 ; O : 3,44 ; Cl : 3,16 ; Br : 2,96. Une liaison est considérée
comme polarisée si $\Delta\chi\geq0{,}4$.
\end{donnees}
\begin{questions}
\item Représenter le schéma de Lewis de chaque molécule.
\item Indiquer les liaisons polarisées et les charges partielles portées par leurs atomes.
\item Pour chaque molécule, identifier un site donneur et un site accepteur de doublet d'électrons.
\item Pourquoi la liaison C--H est-elle considérée comme non polarisée ?
\end{questions}
\end{exercice}
\end{blocLivre}
\begin{blocLivre}{ex-5-3}
\begin{exercice}[1]{Lire des flèches courbes}
À l'état gazeux, l'ammoniac et le chlorure d'hydrogène réagissent en formant des fumées blanches de chlorure
d'ammonium. L'acte élémentaire est représenté ci-dessous, sans ses produits.
\schema{05-ex03}
\begin{donnees}
Électronégativités : H : 2,20 ; N : 3,04 ; Cl : 3,16.
\end{donnees}
\begin{questions}
\item Quelle liaison se forme au cours de cet acte ? Quelle liaison se rompt ?
\item Identifier le site donneur et le site accepteur de doublet d'électrons. Justifier.
\item Écrire les formules des produits, avec leurs charges, et vérifier la conservation de la charge.
\item De quel type de réaction s'agit-il ? Préciser les couples mis en jeu.
\end{questions}
\end{exercice}
\end{blocLivre}
\begin{blocLivre}{ex-5-4}
\begin{exercice}[1]{Catalyseur ou intermédiaire ?}
Au \textsc{xix}\textsuperscript{e} siècle, l'acide sulfurique était fabriqué par le procédé des chambres de plomb, dont
une étape est l'oxydation du dioxyde de soufre, catalysée par le monoxyde d'azote. On la modélise par deux étapes :
\[\text{étape 1 : }\ce{2 NO + O2 -> 2 NO2}\qquad\qquad\text{étape 2 : }\ce{NO2 + SO2 -> NO + SO3}\]
\begin{questions}
\item Établir l'équation de la réaction modélisée par ces deux étapes.
\item Identifier le catalyseur et l'intermédiaire réactionnel. Justifier.
\item Tous les réactifs étant gazeux, la catalyse est-elle homogène ou hétérogène ?
\item L'étape 1 nécessiterait la rencontre simultanée de trois molécules. Pourquoi est-il peu probable qu'il s'agisse
d'un acte élémentaire ?
\end{questions}
\end{exercice}
\end{blocLivre}
\begin{blocLivre}{ex-5-5}
\begin{exercice}[1]{Des chocs plus ou moins efficaces}
L'acte élémentaire \ce{HO- + CH3-Cl -> CH3-OH + Cl-} a lieu lors d'un choc entre un ion hydroxyde et une molécule
de chlorométhane. Les deux situations suivantes représentent un choc efficace et un choc non efficace.
\schema{05-ex05}
\begin{questions}
\item Pourquoi le choc représenté à droite n'est-il pas efficace ?
\item Citer une autre condition nécessaire pour qu'un choc soit efficace.
\item On double la concentration en ions hydroxyde. Comment évolue la fréquence des chocs ? la vitesse de l'acte ?
\item On refroidit le mélange. Expliquer pourquoi la transformation ralentit.
\end{questions}
\end{exercice}
\end{blocLivre}
\rubrique{Entraînement}
\begin{blocLivre}{ex-5-6}
\begin{exercice}[2]{Flèches courbes de réactions acide-base}
On considère les deux actes élémentaires suivants, dont les schémas de Lewis sont donnés.
\schema{05-ex06}
\begin{donnees}
Électronégativités : H : 2,20 ; C : 2,55 ; N : 3,04 ; O : 3,44.
\end{donnees}
\begin{questions}
\item Préciser les couples acide-base mis en jeu dans chaque acte.
\item Identifier, pour chaque acte, le site donneur et le site accepteur de doublet d'électrons.
\item Reproduire les schémas et représenter les flèches courbes de chaque acte.
\item Vérifier la conservation de la charge électrique dans chaque acte.
\end{questions}
\end{exercice}
\end{blocLivre}
\begin{blocLivre}{ex-5-7}
\begin{exercice}[2]{Un acte ou deux ?}
Les deux mécanismes suivants modélisent le remplacement d'un atome de brome par un groupe hydroxyle.
\schema{05-ex07}
\begin{donnees}
Électronégativités : C : 2,55 ; Br : 2,96 ; O : 3,44.
\end{donnees}
\begin{questions}
\item Pour le mécanisme A, identifier le site donneur et le site accepteur. Représenter les flèches courbes.
\item Pour le mécanisme B, identifier l'intermédiaire réactionnel. Quelle particularité présente son atome de carbone central ?
\item Représenter les flèches courbes des deux actes du mécanisme B.
\item Établir l'équation de la réaction modélisée par le mécanisme B.
\end{questions}
\end{exercice}
\end{blocLivre}
\begin{blocLivre}{ex-5-8}
\begin{exercice}[2]{Addition sur une double liaison}
Le propène réagit avec le bromure d'hydrogène selon le mécanisme en deux actes ci-dessous.
\schema{05-ex08}
\begin{donnees}
Électronégativités : H : 2,20 ; C : 2,55 ; Br : 2,96.
\end{donnees}
\begin{questions}
\item Pourquoi la double liaison C=C est-elle un site donneur de doublet d'électrons ? Quel est le site accepteur de
la molécule de bromure d'hydrogène ? Justifier.
\item Représenter les flèches courbes des deux actes.
\item Identifier l'intermédiaire réactionnel.
\item Écrire l'équation de la réaction et nommer le produit formé.
\end{questions}
\end{exercice}
\end{blocLivre}
\begin{blocLivre}{ex-5-9}
\begin{exercice}[2]{Fabriquer de l'éthanol}
L'essentiel de l'éthanol de synthèse est produit par hydratation de l'éthène, en phase gazeuse vers 300 °C, sous 60 à
70 bar, en présence d'acide phosphorique déposé sur un support solide. En solution aqueuse acide, la réaction est
modélisée par le mécanisme suivant.
\begin{center}\renewcommand{\arraystretch}{1.3}
\begin{tabular}{@{}ll@{}}
acte 1 : & \ce{CH2=CH2 + H3O+ -> CH3-CH2+ + H2O}\\
acte 2 : & \ce{CH3-CH2+ + H2O -> CH3-CH2-OH2+}\\
acte 3 : & \ce{CH3-CH2-OH2+ + H2O -> CH3-CH2-OH + H3O+}
\end{tabular}\end{center}
\begin{questions}
\item Établir l'équation de la réaction.
\item Identifier les intermédiaires réactionnels et le catalyseur.
\item Représenter les flèches courbes des actes 1 et 2, à l'aide de schémas de Lewis.
\item Quel est le rôle de l'acte 3 ? Pourquoi l'industrie opère-t-elle vers 300 °C plutôt qu'à température ambiante ?
\end{questions}
\end{exercice}
\end{blocLivre}
\begin{blocLivre}{ex-5-10}
\begin{exercice}[2]{Une catalyse par les ions du fer}
La réaction \ce{S2O8^2- + 2 I- -> 2 SO4^2- + I2} est lente. Elle est catalysée par les ions fer(II), selon le
mécanisme suivant.
\[\text{étape 1 : }\ce{S2O8^2- + 2 Fe^2+ -> 2 SO4^2- + 2 Fe^3+}\qquad\text{étape 2 : }\ce{2 Fe^3+ + 2 I- -> 2 Fe^2+ + I2}\]
On mesure le temps de demi-réaction d'un même mélange dans trois conditions.
\begin{center}\small\renewcommand{\arraystretch}{1.2}
\begin{tabular}{@{}lccc@{}}\toprule
Ajout & aucun & ions \ce{Fe^2+} (1,0 mmol$\cdot$L$^{-1}$) & ions \ce{Fe^3+} (1,0 mmol$\cdot$L$^{-1}$)\\\midrule
$t_{1/2}$ (min) & 40 & 2 & 2\\\bottomrule
\end{tabular}\end{center}
\begin{questions}
\item Vérifier que le mécanisme rend compte de l'équation de la réaction.
\item Quel est le catalyseur ? Pourquoi les ions fer(III) conviennent-ils aussi ?
\item Interpréter, en termes de chocs entre entités, la lenteur de la réaction non catalysée et l'efficacité du catalyseur.
\item Estimer par quel facteur la vitesse de la transformation est multipliée par le catalyseur.
\end{questions}
\end{exercice}
\end{blocLivre}
\begin{blocLivre}{ex-5-11}
\begin{exercice}[2]{Chocs et température}
\begin{docu}{1}{Un modèle des chocs}
La fréquence des chocs entre deux entités est proportionnelle au produit de leurs concentrations et à $\sqrt{T}$, où
$T$ est la température absolue, en kelvins. Expérimentalement, la vitesse de nombreuses réactions double environ
quand la température passe de 20 °C à 30 °C.
\end{docu}
\begin{questions}
\item Par quel facteur la fréquence des chocs est-elle multipliée quand la température passe de 20 °C à 30 °C ?
\item Ce seul effet peut-il expliquer le doublement de la vitesse ? Qu'en déduire sur la proportion de chocs efficaces ?
\item À température constante, on double les concentrations des deux réactifs d'un acte élémentaire mettant en jeu
deux entités. Par quel facteur la vitesse de cet acte est-elle multipliée ?
\item Lors d'une trempe, un prélèvement est dilué cinq fois dans de l'eau glacée. Par quel facteur la fréquence des
chocs entre deux entités réactives est-elle divisée, sans même tenir compte du refroidissement ? Conclure.
\end{questions}
\end{exercice}
\end{blocLivre}
\rubrique{Objectif bac}
\begin{blocLivre}{ex-5-12}
\begin{exercice}[3]{Le chlore et la couche d'ozone}
\begin{docu}{1}{Des atomes de chlore dans la stratosphère}
Les chlorofluorocarbures (CFC), comme le dichlorodifluorométhane \ce{CF2Cl2}, longtemps utilisés dans les
réfrigérateurs et les aérosols, sont très stables. Parvenus dans la stratosphère, ils sont dissociés par le
rayonnement ultraviolet : \ce{CF2Cl2 -> CF2Cl + Cl}. Les atomes de chlore détruisent alors l'ozone selon le mécanisme :
\[\text{acte 1 : }\ce{Cl + O3 -> ClO + O2}\qquad\qquad\text{acte 2 : }\ce{ClO + O -> Cl + O2},\]
les atomes d'oxygène provenant de la dissociation, par le rayonnement, de molécules de dioxygène ou d'ozone.
\end{docu}
\begin{docu}{2}{Un effet démultiplié}
« Un seul atome de chlore peut détruire plus de 100\,000 molécules d'ozone avant d'être éliminé de la stratosphère. »
Signé en 1987, le protocole de Montréal a organisé l'arrêt progressif de la production des CFC.
\source{Agence américaine de protection de l'environnement (EPA), « Basic Ozone Layer Science ».}
\end{docu}
\begin{donnees}
$M(\ce{CF2Cl2})=\qty{120,9}{g.mol^{-1}}$ ; $M(\ce{O3})=\qty{48,0}{g.mol^{-1}}$.
\end{donnees}
\begin{questions}
\item Identifier le catalyseur et l'intermédiaire réactionnel du mécanisme. Établir l'équation de la réaction.
\item Pourquoi une faible quantité de CFC suffit-elle à détruire une grande quantité d'ozone ?
\item \initiative Estimer la masse d'ozone que peuvent détruire les atomes de chlore contenus dans 1,0 kg de
\ce{CF2Cl2}, en admettant qu'ils sont tous libérés et que chacun détruit $\num{1.0e5}$ molécules d'ozone.
\item Proposer une explication au fait que la couche d'ozone ne se reconstitue que lentement, plusieurs décennies
après l'interdiction des CFC.
\end{questions}
\end{exercice}
\end{blocLivre}
\begin{blocLivre}{ex-5-13}
\begin{exercice}[3]{L'arôme de banane}
\begin{docu}{1}{Une synthèse au laboratoire}
L'éthanoate de 3-méthylbutyle, à l'odeur de banane, est synthétisé en chauffant à reflux de l'acide éthanoïque et du
3-méthylbutan-1-ol, noté \ce{R-OH}, en présence de quelques gouttes d'acide sulfurique concentré, qui apporte des
ions \ce{H+}.
\end{docu}
\begin{docu}{2}{Mécanisme simplifié}
\renewcommand{\arraystretch}{1.25}
\begin{tabular}{@{}ll@{}}
acte 1 : & \ce{CH3-COOH + H+ -> CH3-C^+(OH)2}\\
acte 2 : & \ce{CH3-C^+(OH)2 + R-OH -> CH3-C(OH)2-O^+H-R}\\
acte 3 : & \ce{CH3-C(OH)2-O^+H-R -> CH3-C(OH)(O^+H2)-O-R}\\
acte 4 : & \ce{CH3-C(OH)(O^+H2)-O-R -> CH3-C^+(OH)-O-R + H2O}\\
acte 5 : & \ce{CH3-C^+(OH)-O-R -> CH3-COO-R + H+}
\end{tabular}
\end{docu}
\begin{questions}
\item Établir l'équation de la réaction. Identifier le catalyseur et les intermédiaires réactionnels.
\item Dans l'acte 2, identifier le site donneur et le site accepteur de doublet d'électrons, puis représenter la
flèche courbe à l'aide de schémas de Lewis.
\item Dans l'acide éthanoïque, l'atome de carbone du groupe carboxyle porte une charge partielle $\delta^+$. Expliquer
pourquoi l'acte 1 rend la réaction avec l'alcool plus rapide.
\item Pourquoi chauffe-t-on le mélange ? Le chauffage et le catalyseur modifient-ils la quantité d'ester obtenue
lorsque la transformation est terminée ?
\end{questions}
\end{exercice}
\end{blocLivre}
\begin{blocLivre}{ex-5-14}
\begin{exercice}[3]{Le pot catalytique}
\begin{docu}{1}{Dépolluer les gaz d'échappement}
Le pot catalytique « trois voies » d'une voiture à essence contient une céramique en nid d'abeilles, percée de milliers
de canaux et recouverte d'une fine couche contenant quelques grammes de platine, de palladium et de rhodium. Il
transforme le monoxyde de carbone, les hydrocarbures imbrûlés et les oxydes d'azote en dioxyde de carbone, en eau et
en diazote. Il n'est efficace qu'au-delà d'environ 250 à 300 °C. Le plomb, autrefois ajouté à l'essence, se fixe
durablement sur le métal : l'essence plombée a été interdite en France en 2000.
\end{docu}
\schema{05-ex14}
\begin{docu}{2}{Mécanisme en surface}
On note $\ast$ un site libre de la surface du métal et $\mathrm{X}^{\ast}$ une entité X fixée sur un site.
\par\smallskip
{\renewcommand{\arraystretch}{1.25}%
\begin{tabular}{@{}l@{\hspace{10mm}}l@{}}
(1) $\mathrm{CO}+\ast\longrightarrow\mathrm{CO}^{\ast}$ & (4) $\mathrm{CO}^{\ast}+\mathrm{O}^{\ast}\longrightarrow\mathrm{CO_2}+2\,\ast$\\
(2) $\mathrm{NO}+\ast\longrightarrow\mathrm{NO}^{\ast}$ & (5) $\mathrm{N}^{\ast}+\mathrm{N}^{\ast}\longrightarrow\mathrm{N_2}+2\,\ast$\\
(3) $\mathrm{NO}^{\ast}+\ast\longrightarrow\mathrm{N}^{\ast}+\mathrm{O}^{\ast}$ &
\end{tabular}}
\end{docu}
\begin{questions}
\item De quel type de catalyse s'agit-il ? Pourquoi le métal est-il déposé sur une céramique en nid d'abeilles ?
\item \initiative Combiner les étapes du mécanisme pour retrouver l'équation de la réaction :
\[\ce{2 CO + 2 NO -> 2 CO2 + N2}.\]
Identifier les intermédiaires réactionnels. En quoi les sites libres jouent-ils le rôle de catalyseur ?
\item Pourquoi un trajet court, moteur froid, est-il particulièrement polluant ?
\item Expliquer pourquoi le plomb rend le pot catalytique inefficace.
\end{questions}
\end{exercice}
\end{blocLivre}
\begin{blocLivre}{ex-5-15}
\begin{exercice}[3]{La cinétique départage deux mécanismes}
\begin{docu}{1}{Deux mécanismes possibles}
Le remplacement de l'atome de brome d'un bromoalcane RBr par un groupe hydroxyle,
\ce{RBr + HO- -> ROH + Br-}, peut suivre deux mécanismes :
\begin{questions}[label=\otLabel{\Alph*.},leftmargin=1.9em,itemsep=0pt]
\item un seul acte élémentaire : \ce{RBr + HO- -> ROH + Br-} ;
\item deux actes élémentaires : \ce{RBr -> R+ + Br-} (lent), puis \ce{R+ + HO- -> ROH} (rapide).
\end{questions}
\end{docu}
\begin{docu}{2}{Des règles admises}
La vitesse d'une transformation est imposée par son acte élémentaire le plus lent. La vitesse d'un acte élémentaire
est proportionnelle au produit des concentrations des entités qui s'y rencontrent.
\end{docu}
\begin{docu}{3}{Résultats expérimentaux}
Les ions hydroxyde sont toujours en large excès ; on mesure le temps de demi-réaction du bromoalcane.
\par\smallskip
{\small\renewcommand{\arraystretch}{1.2}
\begin{tabular}{@{}lcc@{}}\toprule
 & $[\ce{HO-}]=\qty{0,10}{mol.L^{-1}}$ & $[\ce{HO-}]=\qty{0,20}{mol.L^{-1}}$\\\midrule
2-bromo-2-méthylpropane & 85 s & 85 s\\
bromométhane & 50 min & 25 min\\\bottomrule
\end{tabular}}
\end{docu}
\begin{questions}
\item Pour chaque mécanisme, exprimer la vitesse de la transformation en fonction des concentrations.
\item Attribuer, en justifiant, un mécanisme à chacun des deux bromoalcanes.
\item Pour le 2-bromo-2-méthylpropane, déterminer la constante de vitesse.
\item \initiative Prévoir le temps de demi-réaction du bromométhane si la concentration en ions hydroxyde vaut
\qty{0,050}{mol.L^{-1}}.
\end{questions}
\end{exercice}
\end{blocLivre}
